
\subsection{Sign change of $\eta_k$ and protocol risk}\label{sec:flip}
By \eqref{eq:eta}, all factors of $\eta_k$ other than $x_c-2$ are positive, so the sign of $\eta_k$
equals the sign of $x_c-2$. The function $e^x-1-(k-1)x$ is negative on $(0,x_c)$ and positive on
$(x_c,\infty)$, so it suffices to look at its sign at $x=2$. $e^2-1\approx6.39$, while $(k-1)\cdot2$
equals $6$ for $k=4$ and $8$ for $k=5$. The former gives a positive value, hence $x_c<2$; the latter
gives a negative value, hence $x_c>2$. Since $x_c$ increases with $k$, we have $\eta_k<0$ for
$k=3,4$, while every fixed $k\ge5$ has $\eta_k>0$.
Table~\ref{tab:constants} lists the numerical values.

\begin{table}[!htbp]
\centering
\caption{Analytic constants and the direction of the conditional mean under one-sided failure.}
\begin{tabular}{rrrl}
\toprule
$k$ & $\sigma_k^2$ & $\eta_k$ & direction of the mean ($\rho\ne0$)\\
\midrule
3 & 0.047334402 & $-0.352023542$ & positive\\
4 & 0.022604525 & $-0.041540768$ & positive\\
5 & 0.014632421 & $+0.126662297$ & negative\\
6 & 0.010718241 & $+0.219499658$ & negative\\
\bottomrule
\end{tabular}
\label{tab:constants}
\end{table}

This sign can be translated back to the original scale. Under a one-sided sign, $\rho=1$, and
\eqref{eq:projection} gives $\E[\widehat d\mid F]-d=d(d-1)c_F$, hence $\E[X\mid F]=d(d-1)c_F/s$.
Using $s^2=2d(d-1)/(M-1)$ and $d/M\to\alpha_c$, we obtain
\begin{equation*}
M^{3/2}c_F\longrightarrow
-\frac{\sqrt2\eta_k}{\alpha_c}A(z_r).
\end{equation*}
The sign change is a property of $c_F$ itself inside the critical window, independent of the sign
composition; the sign composition only determines the proportion in which this shift is transmitted
to the estimate. Returning to the interpretation of $c_F$ in Section~\ref{sec:sign}: for $k=3,4$ the
mapping overlap selected by failure is larger than average and the estimate is biased upward; for
$k\ge5$ the overlap is smaller and the estimate is biased downward. For $k=3,4,5,6$, $x_c$ is about
$1.26$, $1.90$, $2.34$, $2.66$. The value of $x_c$ for $k=4$ is close to $2$, so $|\eta_4|$ is only
$0.04$; the selection effect is weakest at $k=4$, and markedly stronger at $k=3$ and for $k\ge5$.
Figure~\ref{fig:selection} gives the simulation comparison for $k=3$--$6$. The mean has another sign
change: for fixed $k\ge5$, as the load decreases gradually the conditional mean changes sign once
more, see Appendix~\ref{app:wm}.

The sign change directly determines on which side the protocol risk falls. Under the rule that takes
$\widehat d/q_u$ as the capacity of the second round, $\widehat d<q_u d$ is equivalent to the
capacity set from the estimate being smaller than the true number of differences, and the lower-tail
probability under failure is the probability of this event. Under a one-sided sign, if $c>0$
($k\ge5$) the true quantile after failure is lower than the quantile of the unconditional normal;
using the latter as the lower-tail threshold makes the lower-tail probability higher than the nominal
value and the capacity too small; if $c<0$ ($k=3,4$) the bias falls on the conservative side. Whether
the second round ultimately succeeds also depends on the capacity conversion factor and the failure
probability of the second-round decoding itself; the end-to-end experiment of Section~\ref{sec:e2e}
checks the coverage of the upper bound after failure, i.e.\ the probability that $U<d$.

The analysis above gives the full conditional distribution when $z_r$ and $c=\eta_k\rho_\infty^2$ are
known. In practice there is a further difficulty: both $z_r$ and $c$ are computed from the true $d$,
and $d$ is precisely the quantity the protocol is meant to estimate. The centering term of the energy
uses only the total sign $b=k^{-1}\sum_vC_v=d\rho$, which the protocol can compute directly.
Section~\ref{sec:upper} uses the observable $(\widehat d,b)$ to perform candidate testing and
inversion, removing the dependence on the true $d$.
